% =========================================================
% Constructing Rational Numbers as Equivalence Classes
% =========================================================

\label{sec:constructing-rational-numbers-as-equivalence-classes}

\subsection*{Definitions and Theorems}

\begin{exposition}
The rational numbers are built from the integers in essentially the same way that the integers were
built from the natural numbers: one begins with concrete representatives, identifies those that ought
to represent the same object, and then passes to equivalence classes.
\end{exposition}

\begin{definition}[Fraction pairs]
\label{def:fraction-pairs}

Let $\mathbb{Z}$ denote the set of integers and define
\[
W = \mathbb{Z} \times (\mathbb{Z}\setminus\{0\}).
\]
Thus
\[
W = \{(a,b) : a,b\in\mathbb{Z},\ b\neq 0\}.
\]

Each ordered pair $(a,b)$ is intended to represent the formal fraction $a/b$.

\end{definition}

\begin{remark*}[Numerator and denominator]
\label{rem:numerator-denominator}

If $(a,b)\in W$, then $a$ is called the \emph{numerator} and $b$ is called the
\emph{denominator}.  At this stage $(a,b)$ is still only an ordered pair; it is not yet a rational
number.

\end{remark*}

\begin{remark*}
Different pairs should represent the same rational number.  For example,
\[
(1,2),\quad (2,4),\quad (3,6)
\]
all ought to describe the same quantity.  To make this precise we introduce an equivalence relation.
\end{remark*}

\begin{definition}[Equality of fraction pairs]
\label{def:fraction-pair-equality}

For $(a,b),(c,d)\in W$, define
\[
(a,b) \sim (c,d)
\quad\Longleftrightarrow\quad
ad = bc.
\]

In fully quantified form,
\[
\forall (a,b),(c,d)\in W,
\qquad
(a,b)\sim(c,d)
\Longleftrightarrow
ad=bc.
\]

\end{definition}

\begin{remark*}
This is the familiar cross-multiplication test for equality of fractions.  For example,
\[
(1,2)\sim(2,4)
\]
because
\[
1\cdot 4 = 2\cdot 2.
\]
\end{remark*}

\begin{theorem}[The Fraction-Pair Relation Is an Equivalence Relation]
\label{thm:fraction-pair-equivalence}
The relation $\sim$ is an equivalence relation on $W$.
\end{theorem}

\begin{remark*}[Proof]
\hyperref[prf:fraction-pair-equivalence]{Go to proof.}
\end{remark*}

\begin{remark*}
Once this is known, the equivalence classes of $\sim$ partition $W$.  Each class is what we shall
call a rational number.
\end{remark*}

\begin{definition}[Rational numbers]
\label{def:rational-numbers}

The set of \emph{rational numbers} is the quotient
\[
\mathbb{Q} = W/\sim.
\]
If $(a,b)\in W$, its equivalence class is denoted by
\[
[(a,b)].
\]
We interpret $[(a,b)]$ as the rational number represented informally by $a/b$.

\end{definition}

\begin{definition}[Distinguished elements]
\label{def:rational-distinguished}

Define
\[
0 := [(0,1)]
\qquad\text{and}\qquad
1 := [(1,1)].
\]

\end{definition}

\begin{remark*}
These definitions are well-defined because
\[
(0,1)\sim(0,b) \quad (b\neq 0)
\qquad\text{and}\qquad
(1,1)\sim(a,a) \quad (a\neq 0).
\]
\end{remark*}
